\answer{ex:11:1}{(a)~$T=k^2/\bigl(\Phi(\Delta I/I)^2\bigr)=9\,\text{s}$. (b)~$45$ sensores, una mancha de $\approx7\times7$; la resolución empeora $\approx7$ veces.}
\answer{ex:11:2}{(a)~$2$--$3.3$ bits por espiga, el $22$--$34\,\%$ del límite ($9.1$--$9.8$ bits). (b)~$\log_2\binom{400}{20}\approx111$ bits; en promedio, $1$ tipo común.}
\answer{ex:11:3}{(a)~$\sigma/9\le I\le9\sigma$, $1.9$ órdenes. (b)~$6$ y $7$.}
\answer{ex:11:4}{(a)~$f<343/0.44\approx780$~Hz. (b)~$625$ espigas, $125$ neuronas.}
\answer{ex:11:5}{$21$ bits por evento, $4.8\cdot10^5$ eventos/s, $7$ eventos por píxel del borde: $136$ píxeles/s. La cámara común, $1.25$ cuadros/s.}
\answer{ex:11:6}{En escala logarítmica, $\approx1.0\,\text{s}$ en ambos sentidos (teoría: $0.96\tau$); en escala lineal, $0.2\,\text{s}$ hacia arriba y $3.0\,\text{s}$ hacia abajo (teoría: $0.18\tau$ y $3.0\tau$).}
\answer{ex:11:7}{(a)~$\ln I$ tiene distribución logística con mediana $\ln\sigma$ y escala $1$: dispersión de $\pi/\sqrt3$ en logaritmos naturales, $0.79$ órdenes. (b)~$r=r_{\max}\Phi_I(I)^2$.}
